\section{Calibração das câmeras}

A calibração final utilizou 22 fotografias da câmera 1x e 20 da câmera 0,5x.
Conforme o histórico experimental informado, a primeira calibração da câmera
0,5x empregou as 22 fotografias; as imagens 019 e 022 apresentaram
os maiores erros individuais e foram retiradas manualmente. Após essa seleção,
o RMS passou de aproximadamente $1{,}686$ para $1{,}463$ pixel. Os arquivos
atuais registram o ajuste final com 20 imagens, sem novos descartes na execução.

Os arquivos \texttt{calibration.npz} e \texttt{log.txt}, nos diretórios de cada
câmera, registram os mesmos parâmetros. As matrizes intrínsecas, em pixels, são
\begin{equation}
K_{1\mathrm{x}}=
\begin{bmatrix}
2715{,}936 & 0 & 1484{,}581\\
0 & 2706{,}305 & 1520{,}002\\
0 & 0 & 1
\end{bmatrix},
\label{eq:k1}
\end{equation}
\begin{equation}
K_{0{,}5\mathrm{x}}=
\begin{bmatrix}
1600{,}726 & 0 & 1586{,}526\\
0 & 1596{,}602 & 1494{,}484\\
0 & 0 & 1
\end{bmatrix}.
\label{eq:k05}
\end{equation}
As menores distâncias focais em pixels da câmera 0,5x são compatíveis com seu
enquadramento mais amplo; não representam medidas físicas em milímetros.

\begin{table}[H]
  \centering
  \caption{Resultados do ajuste final. Os erros são expressos em pixels.}
  \label{tab:calibration}
  \begin{tabular}{lrrr}
    \toprule
    Câmera & Imagens & RMS & Erro médio euclidiano\\
    \midrule
    Principal (1x) & 22 & 0,910 & 0,667\\
    Ultrawide (0,5x) & 20 & 1,463 & 1,168\\
    \bottomrule
  \end{tabular}
\end{table}

A Tabela~\ref{tab:distortion} apresenta os coeficientes adimensionais na
convenção do OpenCV: $k_1$, $k_2$ e $k_3$ são radiais; $p_1$ e $p_2$ são
tangenciais. A notação científica mantém visíveis os coeficientes de menor
magnitude.
\begin{table}[H]
  \centering
  \caption{Coeficientes de distorção das duas câmeras.}
  \label{tab:distortion}
  \begin{tabular}{lrr}
    \toprule
    Coeficiente & Principal (1x) & Ultra-angular (0,5x)\\
    \midrule
    $k_1$ & $1{,}041\times10^{-1}$ & $-1{,}245\times10^{-2}$\\
    $k_2$ & $-3{,}006\times10^{-1}$ & $3{,}508\times10^{-2}$\\
    $k_3$ & $9{,}202\times10^{-5}$ & $1{,}756\times10^{-2}$\\
    $p_1$ & $2{,}460\times10^{-3}$ & $8{,}152\times10^{-4}$\\
    $p_2$ & $1{,}630\times10^{-3}$ & $1{,}142\times10^{-2}$\\
    \bottomrule
  \end{tabular}
\end{table}

O RMS é a raiz da média das distâncias de reprojeção ao quadrado. O erro médio
é a média das distâncias euclidianas entre cantos detectados e reprojetados.
Com 16 pontos por imagem, a média das médias por imagem coincide com a global.

A câmera principal apresentou menor erro nas duas medidas da
Tabela~\ref{tab:calibration}. Campo de visão, tamanho do padrão na imagem,
localização dos cantos e adequação do modelo são possíveis fatores para essa
diferença, mas não foram isolados experimentalmente. Além disso, a redução do
RMS após retirar imagens da câmera 0,5x mede o ajuste ao conjunto selecionado.
